\section{Rangkuman}
\label{sec:rangkuman-bab15}

Pada bab ini telah dibahas:

\begin{itemize}
    \item Array adalah kumpulan data bertipe sama yang disimpan
          berurutan di memori. Array satu dimensi diakses dengan
          register penunjuk dan perulangan.

    \item Array multi dimensi disimpan baris demi baris. Offset
          elemen dihitung dengan
          $(r \times \text{KOLOM} + c) \times \text{ukuran elemen}$.

    \item String adalah array byte yang berisi kode ASCII. Akhir
          string dapat ditandai \code{'\$'} atau dengan menyimpan
          panjangnya.

    \item Penyalinan string dilakukan byte per byte sampai penanda
          akhir. Perbandingan dan pencarian memakai \asm{CMP} di
          dalam perulangan.

    \item Bilangan desimal diubah menjadi digit dengan pembagian
          berulang oleh 10. Sisa pembagian disimpan pada
          \eng{stack} agar urutan cetaknya benar.

    \item Bilangan heksadesimal diubah dengan mengambil empat bit
          sekaligus melalui \asm{ROL} dan \asm{AND} \hex{0F}.

    \item Digit 0--9 diubah menjadi karakter dengan menambahkan
          \hex{30}. Nilai 10--15 diubah menjadi \code{'A'}--\code{'F'}
          dengan menambahkan \hex{37}.

    \item Pencetakan dilakukan dengan layanan DOS
          \code{INT 21h} fungsi \hex{02} untuk satu karakter.
\end{itemize}

\begin{checklistbox}
\vspace{1.1em}
Tandai setiap butir yang sudah Anda kuasai.

\begin{itemize}
    \item[\checkbox] Saya dapat mendeklarasikan array satu dimensi.
    \item[\checkbox] Saya dapat menghitung offset array multi
          dimensi.
    \item[\checkbox] Saya dapat mengakses array dengan \reg{SI} dan
          \asm{LOOP}.
    \item[\checkbox] Saya dapat menyalin, membandingkan, dan mencari
          string.
    \item[\checkbox] Saya dapat mengubah bilangan menjadi digit
          desimal.
    \item[\checkbox] Saya dapat mengubah bilangan menjadi digit
          heksadesimal.
    \item[\checkbox] Saya dapat mencetak hasil konversi dengan
          layanan DOS.
\end{itemize}

\vspace{0.4em}
Bila masih ada butir yang belum bertanda, ulangi bagian yang
bersesuaian sebelum melanjutkan ke bab berikutnya.
\end{checklistbox}

\begin{infobox}
\vspace{1.1em}
\textbf{Rujukan bab.} Materi pada bab ini diringkas dari
\cite{gunawan2019}, Bagian 10, hlm.~93--100, Bagian 14.1,
hlm.~127--131, dan Bagian 16, hlm.~154--158, serta \cite{sto2001},
BAB IX, hlm.~41--46, dan BAB XIX--XX, hlm.~110--121.
\end{infobox}
